\begin{table}[H]
\centering
\caption{Marginal Return to Signal Quality by Delivery Fidelity}
\label{tab:struct_signal_delivery_margins}
\begin{threeparttable}
\scriptsize
\setlength{\tabcolsep}{3pt}
\begin{tabular}[t]{>{\raggedright\arraybackslash}p{3.3cm}ccccc}
\toprule
Scenario & $\omega$ & $\tau$ & Gain $\rho_N\to\rho_K$ & MP of $\rho$ & Cross-partial \\
\midrule
Nigeria realized & 0.84 & 0.31 & $+0.001$ & 0.002 & 0.030 \\
Nigeria high delivery & 0.84 & 0.90 & $+0.088$ & 0.214 & 1.025 \\
Nigeria high delivery + execution & 0.95 & 0.90 & $+0.099$ & 0.241 & 1.157 \\
Kenya observed delivery & 1.00 & 0.50 & $+0.008$ & 0.020 & 0.173 \\
Kenya high delivery & 1.00 & 0.90 & $+0.105$ & 0.254 & 1.218 \\
Fully high-input & 0.98 & 0.95 & $+0.129$ & 0.314 & 1.428 \\
\bottomrule
\end{tabular}
\begin{tablenotes}[para,flushleft]
\footnotesize
\item \textit{Notes:} The table evaluates the assignment-payoff term $\hat{\lambda}\rho\omega\tau^{\hat{\alpha}}$. ``Gain'' is the ATE increase from raising signal quality from Nigeria's estimate to Kenya's estimate, holding the row's $\omega$ and $\tau$ fixed. The last two columns report the analytic marginal product of signal quality and the signal--delivery cross-partial implied by the preferred production map. Country coverage follows the object of the calculation: rows use Nigeria and Kenya because the table varies signal quality between Nigeria's moderate estimate and Kenya's high estimate. Liberia is excluded from this marginal-product table because its role is the low-signal observed cell rather than the benchmark signal-quality upgrade.
\end{tablenotes}
\end{threeparttable}
\end{table}
